\documentclass[12pt, a4paper]{article}

% ── packages ────────────────────────────────────────────────────────────────
\usepackage[T1]{fontenc}
\usepackage[utf8]{inputenc}


\usepackage{amsmath, amssymb, amsthm}
\usepackage{bm}
\usepackage{booktabs}
\usepackage{array}
\usepackage{multirow}
\usepackage{hyperref}
\usepackage{xcolor}
\usepackage{geometry}
\usepackage{natbib}
\usepackage{enumitem}

\geometry{margin=2.5cm}

\hypersetup{
  colorlinks = true,
  linkcolor  = blue!60!black,
  citecolor  = blue!60!black,
  urlcolor   = blue!60!black
}



\title{\textbf{Data Generation Model}\\[4pt]
       \large Spatio-Temporal Extension of \citet{Nalborczyk2025}}
\author{}
\date{\today}

% ── bibliography style ───────────────────────────────────────────────────────
\bibliographystyle{apalike}

% ════════════════════════════════════════════════════════════════════════════
\begin{document}
\maketitle
\tableofcontents
\bigskip

% ════════════════════════════════════════════════════════════════════════════
\section{Data Generation Algorithm}\label{sec:algorithm}

This section describes the complete step-by-step procedure implemented in
\texttt{generate\_data\_2D\_pilot.R}.
The algorithm assembles all components defined in Sections~\ref{sec:signal}--\ref{sec:coords}
into a single simulation loop and produces the long-format dataset of
Section~\ref{sec:output}.

% ────────────────────────────────────────────────────────────────────────────
\subsection{Pre-computation (Outside the Loop)}\label{sec:precompute}

Three quantities are computed once before any loop begins, because they
are shared across all participants and trials.

\paragraph{Step 1 --- Electrode coordinates.}
Raw projected coordinates $(\tilde{x}_c, \tilde{y}_c)$ from the
\texttt{eegcoord} dataset are standardised over all 61 channels
(Equation~\ref{eq:standardise}), then subsetted to the 10 pilot electrodes:
\[
  \hat{x}_c = \tilde{x}_c / \hat{\sigma}_x, \quad
  \hat{y}_c = \tilde{y}_c / \hat{\sigma}_y,
  \quad c \in \mathcal{C}_\text{pilot}
\]

\paragraph{Step 2 --- Spatial signal weights.}
The Gaussian weight of each pilot electrode relative to CZ is computed once:
\[
  w_c = \exp\!\left(-\frac{(\hat{x}_c - \hat{x}_\text{CZ})^2
                          + (\hat{y}_c - \hat{y}_\text{CZ})^2}
                         {2\sigma_s^2}\right), \quad \sigma_s = 1.2
\]
These weights are fixed for the entire simulation.

\paragraph{Step 3 --- Cholesky factor of the spatial noise covariance.}
The inter-electrode distance matrix $D \in \mathbb{R}^{C \times C}$ and
the correlation matrix $\boldsymbol{\Sigma}$ (Equation~\ref{eq:sigma})
are computed once. The lower Cholesky factor $L$ is obtained by:
\[
  \boldsymbol{\Sigma} = LL^\top, \quad
  \boldsymbol{\Sigma}_{cc'} = \exp\!\left(-D_{cc'}^2 / 2\rho^2\right),
  \quad \rho = 0.15
\]
Computing $L$ once and reusing it across all trials is essential for
efficiency: it avoids $R \times K = 30$ redundant factorisations.

% ────────────────────────────────────────────────────────────────────────────
\subsection{Participant Loop}\label{sec:ppt_loop}

For each participant $p = 1, \ldots, P$:

\paragraph{Step 4 --- Sample onset jitter.}
A single onset value is drawn for participant $p$ and held fixed across
all trials:
\[
  \mathrm{onset}_p \sim \mathrm{Uniform}\{150, 152, \ldots, 170\}\,\text{ms}
\]

\paragraph{Step 5 --- Construct ERP template with onset shift.}
The base ERP template $\mathbf{f} \in \mathbb{R}^{N_f}$ (defined in
Section~\ref{sec:erp}, onset fixed at 160\,ms) is time-shifted by
$\delta_p = \mathrm{onset}_p - 160$\,ms through zero-padding:
\[
  \mathbf{f}_p = [0, \ldots, 0,\; f_1, \ldots, f_{93},\; 0, \ldots, 0]
  \in \mathbb{R}^{N_f}, \quad N_f = 251
\]
The number of leading zeros is $\mathrm{onset}_p / 2 - 1$ (indices on the
2\,ms grid), placing the ERP onset precisely at $\mathrm{onset}_p$\,ms.

% ────────────────────────────────────────────────────────────────────────────
\subsection{Trial Loop}\label{sec:trial_loop}

For each trial $r = 1, \ldots, R$ within participant $p$:

\paragraph{Step 6 --- Generate noise for condition 1.}
Two independent noise matrices are drawn:
\begin{align*}
  \mathbf{E}^\text{temporal}_{k=1} &\in \mathbb{R}^{N_f \times C},
    \quad \text{column } c \sim \mathrm{EEGNoise}(\sigma^2/2)
    \text{ independently} \\
  \mathbf{E}^\text{spatial}_{k=1}  &= \mathbf{Z}_{k=1} \cdot L^\top,
    \quad \mathbf{Z}_{k=1} \sim \mathcal{N}(\mathbf{0},\, (\sigma^2/2)\, I_{N_f \times C})
\end{align*}
\[
  \mathbf{N}_{k=1} = \mathbf{E}^\text{temporal}_{k=1}
                   + \mathbf{E}^\text{spatial}_{k=1}
  \in \mathbb{R}^{N_f \times C}
\]

\paragraph{Step 7 --- Generate noise for condition 2.}
An independent draw is made using the same procedure:
\[
  \mathbf{N}_{k=2} = \mathbf{E}^\text{temporal}_{k=2}
                   + \mathbf{E}^\text{spatial}_{k=2},
  \quad \mathbf{N}_{k=2} \perp \mathbf{N}_{k=1}
\]
Condition 1 and condition 2 use \emph{separate} noise draws for every
trial, so the only systematic difference between conditions is the signal.

\paragraph{Step 8 --- Assemble observed EEG.}
For each electrode $c$, the observed time series for the two conditions is:
\begin{equation}\label{eq:assemble}
  y_{c,\cdot,k=1,r} = \mathbf{N}_{k=1}[\,\cdot,\, c]
  \qquad \text{(noise only)}
\end{equation}
\begin{equation}
  y_{c,\cdot,k=2,r} = A \cdot w_c \cdot \mathbf{f}_p
                     + \mathbf{N}_{k=2}[\,\cdot,\, c]
  \qquad \text{(signal + noise)}
\end{equation}
where $A = 1.5$ and $\mathbf{f}_p$ is the participant-specific ERP vector
from Step~5. The signal amplitude scales linearly with $w_c$, so CZ
receives the full signal $A \cdot 1 \cdot \mathbf{f}_p$, while far
electrodes receive $A \cdot w_c \approx 0$.

% ────────────────────────────────────────────────────────────────────────────
\subsection{Post-processing}\label{sec:postproc}

After all loops complete, two post-processing steps are applied.

\paragraph{Step 9 --- Convert to long format.}
All electrode-trial combinations are stacked into a single data frame with
one row per $(p, c, t, k, r)$ observation, and time is converted from
milliseconds to seconds: $t \leftarrow t / 1000$.

\paragraph{Step 10 --- Temporal downsampling.}
Every 5th time point is retained, reducing the sampling rate from
500\,Hz to 100\,Hz (Section~\ref{sec:downsample}).

% ────────────────────────────────────────────────────────────────────────────
\subsection{Sanity Checks}\label{sec:sanity}

Three checks are run after generation to verify correctness.

\paragraph{Signal check.}
For each electrode $c$, the trial-averaged condition difference is computed:
\[
  \overline{\Delta}_c
    = \frac{1}{R}\sum_{r=1}^{R}
      \bigl(\bar{y}_{c,\cdot,k=2,r} - \bar{y}_{c,\cdot,k=1,r}\bigr)
\]
This should be monotonically ordered by $w_c$: CZ largest,
far electrodes (FP1, O1, T7) near zero.

\paragraph{Noise variance check.}
The global standard deviation of \texttt{eeg} across all rows should be
close to $\sqrt{2\,\sigma^2/2} = 1$, since total noise variance is
$\sigma^2 = 1$ per electrode and the signal occupies only a small
fraction of the time series.

\paragraph{Noise correlation check.}
At a single pre-stimulus time point $t < 0$, the observed correlation
between CZ and FCZ under condition 1 is compared to the theoretical value
$\boldsymbol{\Sigma}_\text{CZ,FCZ} = \exp(-D_\text{CZ,FCZ}^2 / 2\rho^2)$.
With only $R = 15$ trials the sample correlation is noisy, but the sign
and rough magnitude should agree.

% ────────────────────────────────────────────────────────────────────────────
\subsection{Algorithm Summary}\label{sec:alg_summary}

\begin{enumerate}[leftmargin=2em, label=\textbf{\arabic*.}]
  \item Standardise electrode coordinates; subset to pilot set.
  \item Compute spatial weights $w_c$ (Gaussian, centred on CZ).
  \item Compute $\boldsymbol{\Sigma}$ and Cholesky factor $L$
        \emph{(once, before any loop)}.
  \item \textbf{For} $p = 1, \ldots, P$:
  \begin{enumerate}[label=\alph*.]
    \item Sample $\mathrm{onset}_p$; build ERP vector $\mathbf{f}_p$.
    \item \textbf{For} $r = 1, \ldots, R$:
    \begin{enumerate}[label=\roman*.]
      \item Draw $\mathbf{N}_{k=1}$ (temporal + spatial noise).
      \item Draw $\mathbf{N}_{k=2}$ (independent).
      \item Assemble $y_{c,t,k,r}$ via Equation~\eqref{eq:assemble}.
    \end{enumerate}
  \end{enumerate}
  \item Stack to long format; convert time to seconds.
  \item Downsample $500\,\text{Hz} \to 100\,\text{Hz}$ (factor 5).
  \item Run sanity checks; save \texttt{sim\_data} and
        \texttt{ground\_truth}.
\end{enumerate}

% ════════════════════════════════════════════════════════════════════════════
\section{Overall Model}\label{sec:overall}

Let $p = 1, \ldots, P$ index participants, $c = 1, \ldots, C$ electrodes,
$t = 1, \ldots, T$ time points, $k \in \{1, 2\}$ conditions, and
$r = 1, \ldots, R$ trials.
The observed EEG value is:
%
\begin{equation}\label{eq:observation}
  y_{p,c,t,k,r}
    = s_{p,c,t} \cdot \mathbf{1}[k = 2]
    + \epsilon_{p,c,t,r}
\end{equation}
%
where the indicator function $\mathbf{1}[k=2]$ is defined as:
\begin{equation}\label{eq:indicator}
  \mathbf{1}[k=2] =
  \begin{cases}
    0 & \text{if } k = 1 \quad \text{(condition 1: noise only)}\\
    1 & \text{if } k = 2 \quad \text{(condition 2: signal + noise)}
  \end{cases}
\end{equation}
The signal $s_{p,c,t}$ and noise $\epsilon_{p,c,t,r}$ are described in
Sections~\ref{sec:signal} and~\ref{sec:noise}, respectively.

% ════════════════════════════════════════════════════════════════════════════
\section{Signal Structure}\label{sec:signal}

The signal $s_{p,c,t}$ is the product of a spatial weight, a temporal
template, and an amplitude scaling factor:
%
\begin{equation}\label{eq:signal}
  s_{p,c,t} = A \cdot w_c \cdot f(t;\,\mathrm{onset}_p)
\end{equation}
%
where $A > 0$ is the signal amplitude parameter that jointly controls
the signal-to-noise ratio together with the noise variance $\sigma^2$
(see Section~\ref{sec:noise}).
In the pilot simulation, $A = 1.5$.
When $A = 1$, Equation~\eqref{eq:signal} reduces to the formulation of
\citet{Nalborczyk2025}.

% ────────────────────────────────────────────────────────────────────────────
\subsection{Temporal Template}\label{sec:erp}

Each participant's onset is drawn from a discrete uniform
distribution:
%
\begin{equation}\label{eq:onset}
  \mathrm{onset}_p
    \sim \mathrm{Uniform}\{150, 152, 154, \ldots, 170\}\ \text{ms}
\end{equation}
%
The step size equals $1000\,\text{ms}/f_s = 2\,\text{ms}$ at $f_s = 500\,\text{Hz},
$so 11 equally probable values are available.
This is a discretised version of the continuous
$\mathrm{Uniform}(150,170)\,\text{ms}$ used in \citet{Nalborczyk2025};
the difference is negligible in practice.

The temporal template $f(t;\,\mathrm{onset}_p)$ is a
\emph{normalised} truncated Gaussian:
%
\begin{equation}\label{eq:erp}
  f(t;\,\mathrm{onset}_p)
  =
  \begin{cases}
    \dfrac{\phi(z_t) - \phi(1.5)}{1 - \phi(1.5)}
    & \text{if } \mathrm{onset}_p \le t \le 342\,\text{ms} \\[8pt]
    0 & \text{otherwise}
  \end{cases}
\end{equation}
%
where $z_t = (t - 250)/50$, $\phi(\cdot)$ is the standard normal
density, $\mu = 250\,\text{ms}$ is the peak, and $\sigma_\text{ERP} = 50\,\text{ms}$
controls the width.
Normalisation ensures $f \in [0,1]$ with a peak value of exactly~$1$ at
$t = 250\,\text{ms}$, so that $A$ in Equation~\eqref{eq:signal} directly
equals the peak signal amplitude.

\paragraph{Implementation note.}
Onset jitter is implemented by time-shifting the fixed base template:
%
\[
  f_p(t) = f(t - \delta_p), \quad
  \delta_p = \mathrm{onset}_p - 160\,\text{ms}
\]
%
i.e.\ the base template (onset fixed at 160\,ms) is shifted by $\delta_p$
milliseconds to the right.
This is numerically equivalent to the truncation definition in
Equation~\eqref{eq:erp} and matches the approach of \citet{Nalborczyk2025}.

% ────────────────────────────────────────────────────────────────────────────
\subsection{Spatial Weights}\label{sec:weights}

EEG signals generated by a cortical source attenuate with distance across
the scalp \citep{NunezSrinivasan2006}.
Let $d_c$ denote the Euclidean distance from electrode $c$ to the activation
centre CZ in standardised coordinate units (see Section~\ref{sec:coords}).
Spatial weights follow a Gaussian decay:
%
\begin{equation}\label{eq:weights}
  w_c
  = \exp\!\left(-\frac{d_c^2}{2\sigma_s^2}\right),
  \quad c = 1, \ldots, C
\end{equation}
%
where $\sigma_s > 0$ controls the spatial spread.
The weight vector satisfies $w_c \in (0,1]$, reaching its maximum
of~$1$ at CZ.

Combining Equations~\eqref{eq:signal} and~\eqref{eq:weights}, the full
signal matrix for participant~$p$ is the outer product:
%
\begin{equation}\label{eq:signal_matrix}
  \mathbf{S}_p = A \cdot \mathbf{w} \otimes \mathbf{f}_p \in \mathbb{R}^{C \times T}
\end{equation}

\paragraph{Choice of activation centre.}
The vertex electrode CZ was selected for two complementary reasons.
Neurophysiologically, the vertex is associated with a broad range of ERP
components, including the P300, the N2, and motor-related potentials, making
it a representative and well-studied location \citep{NunezSrinivasan2006}.
Geometrically, CZ lies at the centre of the 61-electrode layout, so the
Gaussian decay in Equation~\eqref{eq:weights} distributes signal weight
symmetrically across all surrounding electrodes and avoids boundary effects
that would arise from placing the source near the edge of the montage.

% ════════════════════════════════════════════════════════════════════════════
\section{Noise Structure}\label{sec:noise}

The total noise in Equation~\eqref{eq:observation} is the sum of two
independent components:
%
\begin{equation}\label{eq:noise_decomp}
  \epsilon_{p,c,t,r} = \epsilon^{\text{temporal}}_{p,c,t,r}
                      + \epsilon^{\text{spatial}}_{p,c,t,r}
\end{equation}
%
Each component has variance $\sigma^2/2$, so the total single-electrode
noise variance is:
%
\begin{equation}\label{eq:noise_var}
  \mathrm{Var}(\epsilon_{p,c,t,r})
    = \underbrace{\sigma^2/2}_{\text{temporal}}
    + \underbrace{\sigma^2/2 \cdot \Sigma_{cc}}_{= \sigma^2/2,\ \text{spatial}}
    = \sigma^2
\end{equation}
%
In code, \texttt{outvar} $= \sigma^2/2$, not $\sigma^2$.
With \texttt{outvar} $= 0.5$ the total noise variance per electrode
equals $\sigma^2 = 1$.

% ────────────────────────────────────────────────────────────────────────────
\subsection{Temporal Noise}\label{sec:temporal_noise}

Temporal noise mimics EEG spectral characteristics \citep{Yeung2004}.
Formally:
%
\begin{equation}\label{eq:eeg_noise}
  \epsilon^{\text{temporal}}(t_n)
  = \sqrt{\sigma^2/2} \cdot
    \frac{\displaystyle\sum_{i=1}^{50}
          A_i \sin\!\left(\tfrac{2\pi f_i n}{f_s} + \varphi_i\right)
          - \bar{(\cdot)}}
         {\mathrm{sd}(\cdot)}
\end{equation}
%
where the frequency sequence is a cumulative random walk,
$f_i = \sum_{j=1}^{i} 4 u_j$ with $u_j \overset{\text{iid}}{\sim}
\mathrm{Uniform}(0,1)$, covering approximately 1--200\,Hz;
the amplitude weights $A_i = P(f_i)/P(f_1)$ are read from the empirical EEG
power spectrum $P(\cdot)$ \citep{Yeung2004,Rousselet2025};
and $\varphi_i \overset{\text{iid}}{\sim} \mathrm{Uniform}(0, 2\pi)$
are random phases.
The subtraction of the sample mean and division by the sample standard
deviation ensure $\mathrm{Var}(\epsilon^{\text{temporal}}) = \sigma^2/2$
exactly (\texttt{outvar} $= \sigma^2/2$).

Draws are independent across electrodes, trials, and participants,
consistent with \citet{Nalborczyk2025} and \citet{Sassenhagen2019}.

% ────────────────────────────────────────────────────────────────────────────
\subsection{Spatial Noise}\label{sec:spatial_noise}

Neighbouring electrodes share correlated noise due to volume conduction
\citep{NunezSrinivasan2006, Pernet2015}.
The spatial covariance matrix
$\boldsymbol{\Sigma} \in \mathbb{R}^{C \times C}$ is defined as a
\emph{correlation} matrix (diagonal entries equal to~1):
%
\begin{equation}\label{eq:sigma}
  [\boldsymbol{\Sigma}]_{c,c'}
  = \exp\!\left(-\frac{d_{cc'}^2}{2\rho^2}\right)
\end{equation}
%
where $d_{cc'}$ is the Euclidean distance between electrodes $c$ and
$c'$ in \emph{standardised coordinate units} (Section~\ref{sec:coords}),
and $\rho > 0$ is the spatial correlation length.
The noise variance $\sigma^2/2$ is applied via sampling (see below),
so the total spatial covariance is $(\sigma^2/2)\,\boldsymbol{\Sigma}$.

At each time point~$t$, the spatial noise vector across all $C$ electrodes
is drawn as:
%
\begin{equation}\label{eq:spatial_draw}
  \boldsymbol{\epsilon}^{\text{spatial}}_{p,\cdot,t,r}
  = \mathbf{z}_{t,r} \cdot L^\top, \quad
  \mathbf{z}_{t,r} \overset{\text{iid}}{\sim}
    \mathcal{N}\!\left(\mathbf{0},\, \tfrac{\sigma^2}{2} I_C\right)
\end{equation}
%
where $\boldsymbol{\Sigma} = L L^\top$ is the Cholesky factorisation
($L$ lower-triangular).
It follows that
$\mathrm{Cov}(\boldsymbol{\epsilon}^{\text{spatial}}_{p,\cdot,t,r})
= \frac{\sigma^2}{2}\boldsymbol{\Sigma}$,
with diagonal entries $\sigma^2/2$ and off-diagonal entries
$\frac{\sigma^2}{2}\exp(-d_{cc'}^2/2\rho^2)$.

Draws are independent across time points.

% ════════════════════════════════════════════════════════════════════════════
\section{Coordinate Standardisation}\label{sec:coords}

Raw electrode positions are loaded from the \texttt{eegcoord} dataset in the
\texttt{eegkit} R package.
The dataset provides 3D spherical coordinates projected onto a 2D scalp plane
$(\tilde{x}_c, \tilde{y}_c)$ (columns \texttt{xproj}, \texttt{yproj}),
conforming to the extended 10--20 system \citep{NunezSrinivasan2006},
with an approximate physical scale of $1\,\text{unit} \approx 61.5\,\text{mm}$.

Before computing any distances, the projected coordinates are
standardised by their cross-electrode standard deviations:
%
\begin{equation}\label{eq:standardise}
  \hat{x}_c = \frac{\tilde{x}_c}{\hat{\sigma}_x}, \quad
  \hat{y}_c = \frac{\tilde{y}_c}{\hat{\sigma}_y}
\end{equation}
%
where $\hat{\sigma}_x = \mathrm{sd}(\{\tilde{x}_c\}_{c=1}^C)$ and
$\hat{\sigma}_y = \mathrm{sd}(\{\tilde{y}_c\}_{c=1}^C)$ are computed over
all $C = 61$ available electrodes (even when a pilot subset is used for
simulation).

All spatial distances $d_c$ (Eq.~\ref{eq:weights}) and $d_{cc'}$
(Eq.~\ref{eq:sigma}) are computed in standardised units.
The physical interpretation is
$1\,\text{standardised unit} \approx \hat{\sigma}_x \times 61.5\,\text{mm}$.
Consequently, the parameters $\sigma_s$ and $\rho$ are expressed in
standardised units, not in millimetres.

\paragraph{Choice of electrode layout ($C = 61$).}
The 61-channel layout is taken from \texttt{eegcoord} in the \texttt{eegkit}
R package, which provides standardised positions conforming to the extended
10--20 system \citep{NunezSrinivasan2006}.
All 61 available channels are retained without subsetting, so that the
simulated spatial structure reflects a realistic whole-scalp EEG montage.
In the pilot simulation, a 10-electrode subset spanning the full spatial
dynamic range (from CZ to distant controls) is used to reduce computational
cost; the full 61-electrode layout is used in the complete simulation.

% ════════════════════════════════════════════════════════════════════════════
\section{Parameter Settings}\label{sec:params}

\subsection{Fixed Parameters}

Table~\ref{tab:fixed_params} lists all fixed parameters.

\begin{table}[ht]
\centering
\caption{Fixed simulation parameters.
The pilot study uses one participant, 15 trials per condition,
and $T = 51$ time points at 100\,Hz to reduce computational cost.
The full simulation follows \citet{Nalborczyk2025}.}
\label{tab:fixed_params}
\small
\begin{tabular}{llll}
\toprule
\textbf{Parameter} & \textbf{Symbol} & \textbf{Value} & \textbf{Source} \\
\midrule
Participants (pilot / full)
  & $P$ & 1 / 20 & \citealt{Nalborczyk2025} \\
Number of electrodes
  & $C$ & 61 (pilot: 10) & This study \\
Time points (raw / downsampled)
  & $T$ & 251 / 51 & This study \\
Sampling rate (raw / effective)
  & $f_s$ & 500\,Hz / 100\,Hz & This study \\
Downsampling factor
  & --- & 5 & This study \\
Trials per condition (pilot / full)
  & $R$ & 15 / 50 & \citealt{Nalborczyk2025} \\
Number of conditions
  & $K$ & 2 & \citealt{Nalborczyk2025} \\
Total noise variance (per electrode)
  & $\sigma^2$ & 1 & \citealt{Nalborczyk2025} \\
Per-component noise variance
  & $\sigma^2/2$ & 0.5 (\texttt{outvar}) & This study \\
True onset
  & --- & 160\,ms & \citealt{Rousselet2025} \\
True peak
  & --- & 250\,ms & \citealt{Rousselet2025} \\
True offset
  & --- & 342\,ms & \citealt{Rousselet2025} \\
Participant onset range
  & $\mathrm{onset}_p$ & $\{150,152,\ldots,170\}$\,ms & \citealt{Nalborczyk2025} \\
Activation centre
  & --- & CZ & This study \\
Signal amplitude
  & $A$ & 1.5 & This study \\
\bottomrule
\end{tabular}
\end{table}

\subsection{Free Parameters}

Table~\ref{tab:free_params} lists the two free spatial parameters.

\begin{table}[ht]
\centering
\caption{Free parameters introduced by the spatial extension.
All values are in standardised coordinate units
(Section~\ref{sec:coords}).}
\label{tab:free_params}
\small
\begin{tabular}{lllll}
\toprule
\textbf{Parameter} & \textbf{Symbol}
  & \textbf{Meaning}
  & \textbf{Pilot default}
  & \textbf{Sensitivity range} \\
\midrule
Spatial spread
  & $\sigma_s$
  & Breadth of signal (Eq.~\ref{eq:weights})
  & 1.2
  & 0.5, 0.8, 1.2, 1.5 \\
Noise correlation length
  & $\rho$
  & Spatial extent of noise correlation (Eq.~\ref{eq:sigma})
  & 0.15
  & 0.1, 0.15, 0.2, 0.3 \\
\bottomrule
\end{tabular}
\end{table}

In the pilot, $\sigma_s = 1.2$ is set larger than $\rho = 0.15$
to ensure that signal spread and noise correlation have distinguishably
different spatial structures.

% ════════════════════════════════════════════════════════════════════════════
\section{Statistical Models}\label{sec:models}

\subsection{Theoretical background: tensor product smooths}\label{sec:t2_theory}

The spatio-temporal smooth $f(\hat{x}, \hat{y}, t)$ is constructed as a
\emph{tensor product smooth} via the \texttt{t2()} constructor in
\texttt{mgcv}/\texttt{brms}.
This choice is non-trivial and follows the theory established in
\citet{Wood2006tensor, WoodScheiplFaraway2013} and \citet[Ch.~5]{Wood2017}.
\citet{Nalborczyk2025} adopt the same constructor (their Appendix~A and
Section~3.2) and cite \citet{Wood2017} when introducing it.
We summarise here the four theoretical points that justify our concrete
modelling choices.

\subsubsection{Tensor product construction}\label{sec:theory_tensor}

A 1D smooth $g(x) = \sum_{i=1}^{k_x} \alpha_i B_i(x)$ generalises to two
dimensions \emph{not} by inflating it to a single 2D basis, but by allowing
each coefficient $\alpha_i$ to itself vary smoothly along the second
covariate $z$, i.e.\ $\alpha_i(z) = \sum_{j=1}^{k_z} \beta_{ij} C_j(z)$.
The resulting surface is
%
\begin{equation}\label{eq:tensor_basis}
  f(x, z) = \sum_{i=1}^{k_x}\sum_{j=1}^{k_z}
            \beta_{ij}\, B_i(x)\, C_j(z),
\end{equation}
%
a linear combination of the $k_x \times k_z$ products of marginal basis
functions \citep[\S 5.6.1]{Wood2017}.
Geometrically (cf.\ \citet[Fig.~5.17]{Wood2017}), this corresponds to a
\emph{lattice} of marginal strips rather than a single thin-plate sheet;
each marginal direction retains its own stiffness and unit.
Our model uses $d = c(2, 1)$, which groups the two spatial coordinates
$(\hat{x}, \hat{y})$ as a single 2D margin and treats $t$ as a separate 1D
margin.

\subsubsection{Anisotropy and scale invariance}\label{sec:theory_aniso}

Isotropic smooths such as $s(\hat{x}, \hat{y}, t)$ implicitly assume that
displacements along the three covariates are commensurable.
This is reasonable for geographic coordinates measured in the same unit but
fails when, as here, $\hat{x}$ and $\hat{y}$ are in standardised spatial
units while $t$ is in seconds: the implied relative importance of spatial
versus temporal smoothness depends on an arbitrary rescaling of $t$
\citep[\S 5.7]{Wood2017}.
Tensor product smooths resolve this by assigning each marginal direction
an independent smoothing parameter, yielding fits that are invariant under
linear rescaling of any single covariate.

\paragraph{Isotropy within vs.\ across marginals.}
A subtle but important point is that the thin plate basis chosen for each
marginal is itself isotropic \citep[\S 5.7]{Wood2017}: rotations of the
marginal's covariate space leave the penalty invariant. This
within-marginal isotropy is appropriate here because the two spatial
coordinates $(\hat{x}, \hat{y})$ share the same standardised unit
(Section~\ref{sec:coords}), and the temporal marginal is one-dimensional,
for which isotropy is trivially satisfied. The tensor product construction
preserves isotropy \emph{within} each marginal while introducing
anisotropy \emph{across} marginals via the independent smoothing
parameters, thereby reconciling the thin plate basis choice with the
scale incommensurability between space and time.

\subsubsection{Basis dimension and the rôle of the penalty}\label{sec:theory_k}

The basis dimension $k$ in
\texttt{t2(}$\ldots$\texttt{, k = c(}$k_{xy},\,k_t$\texttt{))} sets an
\emph{upper bound} on the dimension of the function space the smooth
can occupy. The \emph{effective} degrees of freedom are not $k$ but are
determined by the smoothing parameters $\boldsymbol{\lambda}$, which
shrink penalised coefficients toward the unpenalised null space
\citep[\S 5.9]{Wood2017}. Increasing $k$ therefore does not by itself
cause overfitting; it only enlarges the search space in which the
penalty operates. The standard practical recommendation is to choose
$k$ large enough that diagnostic checks (\texttt{gam.check()},
\citealp[\S 5.9.2]{Wood2017}; \citealp{PyaWood2016}) indicate the basis
is not the binding constraint.

In the 1D univariate case this guidance is well operationalised:
\citet[Appendix~B]{Nalborczyk2025} sweep $k \in \{5, 10, 20, 40, 80\}$
and select $k=40$ via LOOIC. No closed-form analogue is available for
tensor-product BGAMM in 3D, however. We therefore adopt a
\emph{two-phase protocol}: derive a starting point from each marginal
considered independently (Phase~1), then validate against three
constraints and an empirical sensitivity sweep (Phase~2).

\paragraph{Phase 1 — Marginal-based derivation.}
Following \citet[\S 4]{Pedersen2019}, each marginal's $k$ can be chosen
independently from its own data-driven requirements:
\begin{itemize}
  \item \textbf{Temporal marginal ($k_t$).}
        Choose $k_t$ such that the basis spacing $T_{\text{window}}/k_t$
        is small relative to the narrowest signal feature. For the
        $\sigma_{\text{ERP}} = 50$\,ms truncated Gaussian
        (Section~\ref{sec:erp}) the temporal full-width at half-maximum
        is $\approx 118$\,ms, demanding basis spacing $\lesssim 40$\,ms;
        with a 500\,ms window this implies $k_t \gtrsim 13$. We use
        $k_t \in [20, 30]$, broadly consistent with the $k = 40$ chosen
        by \citet{Nalborczyk2025} in the 1D analysis.
  \item \textbf{Spatial marginal ($k_{xy}$).}
        For a 2D thin plate over $n_{\mathrm{ch}}$ electrodes, the
        usual heuristic is $k_{xy} \sim n_{\mathrm{ch}}$, possibly
        slightly larger \citep{Wood2017}. The pilot uses
        $n_{\mathrm{ch}} = 10$, so $k_{xy} \in [10, 20]$ is appropriate;
        choosing $k_{xy} = 30$ corresponds to three times the spatial
        resolution actually present in the data and risks parameter
        bloat. The full study uses $n_{\mathrm{ch}} = 61$, so
        $k_{xy} \in [25, 35]$ is reasonable.
\end{itemize}
For our pilot, Phase~1 therefore suggests candidate configurations
$(k_{xy}, k_t) \in \{(10, 20), (10, 30), (20, 20), (20, 30)\}$.

\paragraph{Phase 2 — Constraint validation and sweep.}
The Phase~1 candidates are then checked against three constraints:

\textit{C1 --- Null-space lower bound.} The dimension of the unpenalised
null space of an $m$-th order thin plate spline in $d$ dimensions is
$M = \binom{m+d-1}{d}$ \citep[\S 5.5]{Wood2017}. A meaningful penalty
requires $k > M$. For $m=2$: $k_{xy} > 3$, $k_t > 2$, satisfied by all
Phase~1 candidates.

\textit{C2 --- Bayesian information density.} Stable Bayesian inference
in sparse-data regimes requires roughly five or more cells per
coefficient \citep[informal rule discussed in][]{Gelman2006}. With
$N$ design cells, $C$ condition levels, and $B$ ANOVA blocks per
\texttt{by=condition} smooth under \texttt{full=TRUE},
\[
   k_{xy}\,k_t \;\lesssim\; \frac{N}{5\,C\,B}.
\]
For our pilot ($N \approx 5100$, $C=2$, $B=6$) this yields
$k_{xy} k_t \lesssim 85$. All Phase~1 candidates satisfy this. The
exploratory configurations $(30, 20)$ and $(30, 30)$ with
\texttt{full=TRUE} \emph{violate} this constraint
($k_{xy} k_t \in \{600, 900\}$), making them
information-density-borderline or worse.

\textit{C3 --- Diagnostic check.} For each fitted configuration we run
\texttt{gam.check()} \citep[\S 5.9.2]{Wood2017} and verify that the
effective degrees of freedom (EDF) are well below the reference
degrees of freedom $k'$ (we observe $\mathrm{EDF}/k' \approx 0.35$ for
the leading smooth in our pilot fits, confirming the basis is not
binding).

\paragraph{Sensitivity sweep.}
Within the resulting bracket we fitted BGAMM with seven $k$
configurations:
$(k_{xy}, k_t) \in \{(10, 12),\, (10, 20),\, (10, 30),\,
(20, 20),\, (20, 30),\, (30, 20),\, (30, 30)\}$,
covering Phase~1-compliant configurations as well as the
exploratory over-budget cases. Each fit was evaluated on:
peak-amplitude recovery at the seed electrode (Target~4,
Section~\ref{sec:targets}); onset/offset bias (Target~2);
spatial and spatio-temporal $F_1$ in 30-sim Monte Carlo
(Section~\ref{sec:theta_sens}); and \texttt{gam.check} k-index.
Configurations violating C2 ($(30, 20)$ and $(30, 30)$ with
\texttt{full=TRUE}) exhibit pronounced amplitude under-recovery
(36\% and lower) consistent with the curse-of-dimensionality theory
of Section~\ref{sec:amp_theory}. Phase~1-compliant configurations
recover 46--55\% amplitude under sparse pilot data with stable F1.

\paragraph{Recommended default.}
Combining Phase~1 reasoning, the three constraints C1--C3, and the
empirical sweep, our recommended default for the pilot is
$(k_{xy}, k_t) = (20, 20)$: it satisfies all constraints, achieves
the best Monte Carlo $F_1$ score (0.80) and timing bias
($\pm 10$\,ms at threshold~$=10$) across 30 simulations, and aligns
with the marginal-based reasoning above. For the full study
($n_{\mathrm{ch}} = 61$, larger $N$) the constraints loosen and
$(k_{xy}, k_t) \in \{(25, 30), (30, 30)\}$ become viable.

\paragraph{Relation to \citet{Nalborczyk2025}'s 1D protocol.}
\citet[Appendix~B]{Nalborczyk2025} use a two-step procedure for 1D
$k$-selection: (i) an initial value chosen by intuition (``$k=20$
seemed appropriate''), and (ii) a sensitivity sweep over
$k \in \{5, 10, 20, 40, 80\}$ with LOOIC-based comparison. This
informal initialisation is tractable in 1D because the search space
is small. The 3D tensor-product setting introduces two complications:
the search space becomes a Cartesian product $k_{xy} \times k_t$,
and each HMC fit takes approximately three days on our cluster. A
naive 5$\times$5 sweep would require $\sim$25 fits, or $\sim$75
cluster-days, which is infeasible for our project. Phase~1 of our
protocol therefore replaces the informal initialisation with explicit
marginal-based reasoning \citep[\S 4]{Pedersen2019}, anchoring $k_t$
to the signal FWHM and $k_{xy}$ to the number of electrodes; the
result is a focused 4--7-fit sensitivity sweep (Phase~2) that retains
the LOOIC + PPC validation of \citet{Nalborczyk2025} while explicitly
justifying the candidate region.

\subsubsection{ANOVA decomposition (\texttt{full = TRUE})}\label{sec:theory_full}

The default \texttt{t2()} construction (\texttt{full = FALSE}) assigns one
penalty per margin and ties marginal-direction coefficients into a single
nested basis \citep{Wood2006tensor}.
The \texttt{full = TRUE} variant decomposes the smooth into an
ANOVA-style hierarchy
%
\begin{equation}\label{eq:anova}
  f(\hat{x}, \hat{y}, t)
  = f_{xy}(\hat{x}, \hat{y})
  + f_{t}(t)
  + f_{xy,t}(\hat{x}, \hat{y}, t),
\end{equation}
%
where each component carries its own penalty
\citep{WoodScheiplFaraway2013}.
The interaction term $f_{xy,t}$, which encodes precisely the
``where-and-when'' information that distinguishes evoked EEG effects from
purely spatial or purely temporal trends, then receives a penalty
independent of the spatial and temporal main effects, mitigating the risk
that an inappropriately strong main-effect penalty over-shrinks
interaction structure \citep[\S 5.6.3]{Wood2017}.
\citet{Nalborczyk2025} (Section~3.2) explicitly fit their tensor product
``with full penalties'' citing \citet{Wood2017}; we follow this convention.

\subsubsection{\texorpdfstring{Bayesian inference of $(\boldsymbol{\beta}, \boldsymbol{\lambda})$}{Bayesian inference of (beta, lambda)}}\label{sec:theory_inference}

In the penalised-likelihood formulation, the penalty $\lambda
\boldsymbol{\beta}^{\!\top}\!\mathbf{S}\boldsymbol{\beta}$ is equivalent
to a Gaussian prior
$\boldsymbol{\beta}\!\mid\!\boldsymbol{\lambda} \sim
\mathcal{N}(\mathbf{0}, (\lambda\mathbf{S})^{-1})$, and valid inference
requires the \emph{joint} posterior of $(\boldsymbol{\beta},
\boldsymbol{\lambda})$ \citep[\S 6.10--6.11]{Wood2017,MarraWood2012}.
We approximate this joint posterior with Hamiltonian Monte Carlo (NUTS)
in \texttt{brms}/\texttt{Stan}, which preserves the strong correlations
between spline coefficients and smoothing parameters.
Mean-field variational approximations factorise $q(\boldsymbol{\beta},
\boldsymbol{\lambda}) = q(\boldsymbol{\beta})\,q(\boldsymbol{\lambda})$
and are known to systematically over-shrink $\boldsymbol{\beta}$ and
under-estimate posterior variance in this setting; they are therefore
avoided for inferential statements about the recovered surface.

\subsection{Spatio-Temporal BGAM}

The modelling goal is to estimate where and when the two conditions differ.
\citet{Nalborczyk2025} (Appendix~A, p.~26) note that the 1D temporal BGAM
can be extended to spatial and temporal dimensions via
\texttt{te(x, y, time, d = c(2, 1))}; the present study implements this
extension using the \texttt{t2()} constructor for full marginal penalties.

\subsubsection{Single-participant model (pilot)}

For $P = 1$, the model contains two additive smooth components:
%
\begin{equation}\label{eq:model_pilot}
  \mathrm{eeg}_{c,t,k}
  = \alpha
    + f_0(\hat{x}_c, \hat{y}_c, t)
    + f_\Delta(\hat{x}_c, \hat{y}_c, t) \cdot \mathbf{1}[k=2]
    + \varepsilon_{c,t,k}
\end{equation}
%
where $f_0$ is the baseline spatio-temporal smooth (shared across conditions)
and $f_\Delta$ is the condition-difference smooth.
In \texttt{brms} syntax:
%
\begin{verbatim}
eeg ~ condition +
  t2(xproj, yproj, time, bs = c("tp","tp","tp"),
     k = c(10, 10), d = c(2, 1)) +
  t2(xproj, yproj, time, bs = c("tp","tp","tp"),
     k = c(10, 10), d = c(2, 1), by = condition)
\end{verbatim}

\subsubsection{Multi-participant model (full study)}

For $P > 1$, a varying intercept $u_p \sim \mathcal{N}(0,\tau^2)$ accounts
for between-participant variability \citep{Nalborczyk2025}:
%
\begin{equation}\label{eq:model_full}
  \mathrm{eeg}_{p,c,t,k}
  = \alpha
    + f_0(\hat{x}_c, \hat{y}_c, t)
    + f_\Delta(\hat{x}_c, \hat{y}_c, t) \cdot \mathbf{1}[k=2]
    + u_p
    + \varepsilon_{p,c,t,k}
\end{equation}
%
with \texttt{(1|participant)} added to the \texttt{brms} formula.
The argument \texttt{d = c(2, 1)} specifies that $\hat{x}$ and $\hat{y}$
form a 2D spatial block, while \texttt{time} enters as a separate 1D
dimension, following the advice in \citet{Nalborczyk2025} (Appendix~A).

\subsection{Posterior Odds and Cluster Detection}

The quantity of interest is the condition-difference smooth
$f_\Delta(c,t)$.
Following \citet{Nalborczyk2025} (Section~2.1.2, Figure~4), we compute
the posterior probability that the condition effect is positive:
%
\begin{equation}\label{eq:prob}
  p_{c,t} = P\!\left(f_\Delta(\hat{x}_c, \hat{y}_c, t) > 0 \mid \text{data}\right)
\end{equation}
%
estimated from $S$ posterior draws as:
%
\begin{equation}\label{eq:prob_est}
  \hat{p}_{c,t} = \frac{1}{S}\sum_{s=1}^{S}
    \mathbf{1}\!\left[\hat{y}^{(s)}_{c,t,k=2} - \hat{y}^{(s)}_{c,t,k=1} > 0\right]
\end{equation}
%
The posterior odds are:
%
\begin{equation}\label{eq:odds}
  \mathrm{OR}_{c,t} = \frac{\hat{p}_{c,t}}{1 - \hat{p}_{c,t}}
\end{equation}
%
A node $(c,t)$ is classified as active if $\mathrm{OR}_{c,t} > \theta$,
where $\theta$ is the threshold (default $\theta = 10$, following
\citealt{Nalborczyk2025}).
Active nodes that are adjacent in both space (electrode neighbours) and
time form a spatio-temporal cluster.

\subsection{Posterior Odds Threshold Sensitivity}

Following \citet{Nalborczyk2025} (Section~2.1.2), we evaluate
$\theta \in \{3, 5, 10, 20, 50\}$ without refitting the model.

\subsection{Spatio-Temporal Cluster-Based Permutation Test}

As a comparison method, we implement the spatio-temporal cluster-based
permutation test of \citet{Maris2007,Pernet2015}:
%
\begin{enumerate}[label=\arabic*.]
  \item Compute a mass-univariate $t$-statistic at each node $(c,t)$.
  \item Apply a cluster-forming threshold ($p < 0.05$ uncorrected) to
        identify candidate clusters of contiguous active nodes.
  \item Compute the cluster mass (sum of squared $t$-values) per cluster.
  \item Construct a null distribution of maximum cluster masses via
        1{,}000 random permutations of condition labels.
  \item Declare a cluster significant if its mass exceeds the 95th percentile
        of the null distribution.
\end{enumerate}

% ════════════════════════════════════════════════════════════════════════════
\section{Data Output Format}\label{sec:output}

\subsection{Structure}

Simulated data are stored in long format (Table~\ref{tab:output_cols}).

\begin{table}[ht]
\centering
\caption{Long-format data structure.}
\label{tab:output_cols}
\small
\begin{tabular}{llll}
\toprule
\textbf{Variable} & \textbf{Type} & \textbf{Description} & \textbf{Example} \\
\midrule
\texttt{participant} & character & Participant ID & \texttt{participant\_01} \\
\texttt{channel}     & character & Electrode label & \texttt{CZ} \\
\texttt{xproj}       & numeric   & Standardised $x$ coordinate & $-0.23$ \\
\texttt{yproj}       & numeric   & Standardised $y$ coordinate & $0.14$ \\
\texttt{condition}   & factor    & Experimental condition & \texttt{cond1}, \texttt{cond2} \\
\texttt{trial}       & integer   & Trial index & $1, 2, \ldots, 15$ \\
\texttt{time}        & numeric   & Time in seconds & $0.000, 0.010, \ldots, 0.500$ \\
\texttt{eeg}         & numeric   & EEG amplitude & $0.32, -0.14, \ldots$ \\
\bottomrule
\end{tabular}
\end{table}

Total rows for the pilot ($P=1$, $C=10$, $T=51$, $K=2$, $R=15$):
$1 \times 10 \times 51 \times 2 \times 15 = 15{,}300$.\\
Total rows for the full simulation ($P=20$, $C=61$, $T=250$, $K=2$, $R=50$):
$20 \times 61 \times 250 \times 2 \times 50 = 30{,}500{,}000$.

\subsection{Temporal Downsampling}\label{sec:downsample}

After generation at 500\,Hz, every $k$-th sample is retained
(naive downsampling, no anti-aliasing filter applied).
The pilot uses $k = 5$, yielding:
%
\begin{equation}
  \mathcal{T}_\text{keep}
    = \{t_n \in \mathcal{T}_\text{sec} : n \equiv 1\ (\mathrm{mod}\ 5)\}
    = \{0.000, 0.010, 0.020, \ldots, 0.500\}\,\text{s}
\end{equation}
%
giving $\lfloor(251-1)/5\rfloor + 1 = 51$ time points at
$f_s^\text{eff} = 100\,\text{Hz}$.

\textbf{Limitation (D14).} No low-pass filter precedes downsampling.
EEG noise covers frequencies up to $\approx$\,200\,Hz; aliasing may
contaminate the noise spectrum above 50\,Hz.
This has minimal impact on ERP detection (signal energy concentrated at
$<$\,30\,Hz) but should be addressed in the full simulation via a
50\,Hz Butterworth low-pass filter before subsampling.

\subsection{Comparison with the Original Simulation}

Table~\ref{tab:comparison} compares this study with \citet{Nalborczyk2025}.

\begin{table}[ht]
\centering
\caption{Comparison between the original and extended simulations.}
\label{tab:comparison}
\small
\begin{tabular}{lll}
\toprule
 & \textbf{Nalborczyk \& Bürkner (2025)} & \textbf{This study} \\
\midrule
Signal structure
  & $f(t;\,\mathrm{onset}_p)$
  & $A \cdot w_c \cdot f(t;\,\mathrm{onset}_p)$ \\
Noise structure
  & $\epsilon^\text{temporal}$ only
  & $\epsilon^\text{temporal} + \epsilon^\text{spatial}$ \\
Electrode independence
  & Electrodes independent
  & Inter-electrode correlation via $\boldsymbol{\Sigma}$ \\
Model formula
  & \texttt{s(time) + (1|id)}
  & \texttt{t2(x,y,t,d=c(2,1), by=cond)} \\
Comparison method
  & Cluster permutation (1D)
  & Cluster permutation (2D) \\
Free parameters
  & None
  & $\sigma_s,\, \rho,\, A,\, \theta$ \\
\bottomrule
\end{tabular}
\end{table}

% ════════════════════════════════════════════════════════════════════════════
\section{Ground Truth Definition}\label{sec:ground_truth}

The ground truth defines the true spatio-temporal activation region
$\mathcal{A}$, against which cluster detection is evaluated.

\paragraph{Temporal dimension.}
The true activation window is $[\text{onset}, \text{offset}] = [0.160, 0.342]\,\text{s}$.
After downsampling, the active time points are:
%
\[
  \mathcal{T}^\text{ds}_\text{active}
    = \{t \in \mathcal{T}_\text{keep} : 0.160 \le t \le 0.340\}
\]
%
Note: $0.342\,\text{s}$ does not fall on a 100\,Hz grid point
($0.342 / 0.010 = 34.2$).
The last active grid point is $0.340\,\text{s}$.
Ground truth stores \texttt{true\_offset\_ds = 0.340}.

\paragraph{Spatial dimension.}
Active electrodes are defined by a signal-to-noise threshold:
%
\begin{equation}
  \mathcal{C}_\text{active}
    = \left\{c : A \cdot w_c > \sqrt{\sigma^2}\right\}
    = \left\{c : w_c > \frac{1}{A} = \frac{1}{1.5} \approx 0.667\right\}
\end{equation}
%
This threshold is grounded in signal-to-noise theory: an electrode
is considered ``active'' if its peak signal exceeds the total noise standard
deviation.
Code uses \texttt{spatial\_weight > 0.1} as a provisional display threshold
(which is more permissive); the stricter $w_c > 0.667$ is recommended for
formal evaluation of detection performance.

\paragraph{True activation set.}
%
\[
  \mathcal{A}
    = \mathcal{C}_\text{active} \times \mathcal{T}^\text{ds}_\text{active}
    \subset \{1,\ldots,C\} \times \{1,\ldots,T\}
\]

Detection performance metrics are then:
%
\[
  \text{Sensitivity}
    = \frac{|\{(c,t) \in \mathcal{A} : \widehat{\mathrm{OR}}_{c,t} > \theta\}|}
           {|\mathcal{A}|},
  \quad
  \text{Specificity}
    = \frac{|\{(c,t) \notin \mathcal{A} : \widehat{\mathrm{OR}}_{c,t} \le \theta\}|}
           {|\mathcal{A}^c|}
\]

% ════════════════════════════════════════════════════════════════════════════
\section{Connection to Nalborczyk \& Bürkner (2025)}\label{sec:connection}

Table~\ref{tab:mapping} maps each equation in the present document to the
corresponding definition or figure in \citet{Nalborczyk2025}.

\begin{table}[ht]
\centering
\caption{Cross-reference between this document and \citet{Nalborczyk2025}.
All page numbers refer to the bioRxiv preprint (August 2025).}
\label{tab:mapping}
\small
\begin{tabular}{p{5cm} p{4.5cm} p{4.5cm}}
\toprule
\textbf{This document} & \textbf{Nalborczyk \& Bürkner (2025)} & \textbf{Notes} \\
\midrule
Eq.~\ref{eq:observation}: observation model
  & Section 2.1.1, p.~8
  & Original is 1D; spatial index $c$ is new \\
Eq.~\ref{eq:erp}: ERP template (normalised)
  & Section 2.1.1, p.~8; \texttt{erp\_template.R}
  & Nalborczyk uses same code; normalisation implicit \\
Eq.~\ref{eq:eeg_noise}: EEGNoise
  & Section 2.1.1, p.~8; \texttt{eeg\_noise.R}
  & Identical function; variance control via \texttt{outvar} \\
Eq.~\ref{eq:prob}--\ref{eq:odds}: posterior odds
  & Section 2.1.2, Figure~4, p.~10
  & Identical formula; extended from 1D to $(c,t)$ nodes \\
$\theta = 10$ threshold
  & Figure~4 caption, p.~10
  & Default value identical \\
\texttt{t2(x,y,time,d=c(2,1))}
  & Appendix~A, p.~26
  & Nalborczyk mentions this exact extension \\
Varying intercept $u_p$
  & Section 2.1.2, p.~9
  & Identical; from \citet{Nalborczyk2019} \\
Sensitivity range $\theta \in \{3,5,10,20,50\}$
  & Section 2.1.3, p.~10; Nalborczyk pers.\ comm.
  & Identical \\
Cluster-based permutation (2D)
  & Section 2.1.3, p.~10
  & Original is 1D; spatial neighbours added \\
\bottomrule
\end{tabular}
\end{table}

% ════════════════════════════════════════════════════════════════════════════
\section{Model Fitting Pipeline (\texttt{01\_fit\_model.R})}\label{sec:fitting}

This section describes the steps taken to prepare data, fit the BGAM, and
validate the fit via three diagnostic targets.
The pipeline corresponds to \texttt{01\_fit\_model.R} and follows the
recommendations of \citet{Nalborczyk2025} (Sections 2.1.2--2.1.4).

% ────────────────────────────────────────────────────────────────────────────
\subsection{Data Reduction}\label{sec:reduction}

Starting from the pilot dataset
($C = 10$, $T = 51$, $K = 2$, $R = 15$),
two further reductions are applied before model fitting.

\paragraph{Trial truncation.}
Only the first $R' = \min(R,\,20)$ trials per condition are retained.
In the pilot $R = 15 < 20$, so this step has no effect; it serves as a
safety cap for the full simulation ($R = 50$).

\paragraph{Temporal sub-sampling.}
Every other time point is kept:
\begin{equation}\label{eq:subsample}
  \mathcal{T}^* = \bigl\{t_n \in \mathcal{T}'' : n \text{ odd}\bigr\},
  \qquad |\mathcal{T}^*| = 26, \quad f_s^* \approx 50\,\text{Hz}
\end{equation}
Combined with the earlier 500\,Hz $\to$ 100\,Hz step, the data entering
the model have been downsampled by a total factor of~10.

\paragraph{Trial averaging.}
For each node $(c, t, k)$, EEG values are averaged across retained trials:
\begin{equation}\label{eq:trial_avg}
  \bar{y}_{c,t,k} = \frac{1}{R'}\sum_{r=1}^{R'} y_{c,t,k,r}
\end{equation}
This is the ``computationally lighter version fitted on
by-participant summary statistics'' recommended by \citet{Nalborczyk2025}
(Section 2.1.2, p.\,9).
Averaging reduces the per-electrode noise variance from $\sigma^2 = 1$
to $\sigma^2 / R'$, improving SNR by a factor of $\sqrt{R'}$.

% ────────────────────────────────────────────────────────────────────────────
\subsection{Statistical Model}\label{sec:brms_model}

The trial-averaged data $\bar{y}_{c,t,k}$ are modelled as:
\begin{equation}\label{eq:brms}
  \bar{y}_{c,t,k}
    = \alpha + \beta_k
    + f_{\text{cond1}}(\hat{x}_c, \hat{y}_c,\, t)
    + f_{\text{cond2}}(\hat{x}_c, \hat{y}_c,\, t)
    + \varepsilon_{c,t,k},
  \quad \varepsilon_{c,t,k} \sim \mathcal{N}(0, \sigma_\varepsilon^2)
\end{equation}
where $\alpha$ is the global intercept, $\beta_k$ a condition main effect,
and $f_{\text{cond1}}, f_{\text{cond2}}$ are condition-specific
spatio-temporal smooths estimated via \texttt{by = condition}.
The condition-difference smooth is $f_\Delta = f_{\text{cond2}} - f_{\text{cond1}}$,
matching Equation~\eqref{eq:model_pilot}.

Each smooth is a tensor-product thin-plate spline
(\citealt{Nalborczyk2025}, Appendix~A):
\begin{equation}\label{eq:basis}
  f_k(\hat{x}_c, \hat{y}_c, t)
    = \sum_{j=1}^{K_s}\sum_{l=1}^{K_t}
      \beta_{jl}^{(k)}\,
      B_j^{(s)}(\hat{x}_c, \hat{y}_c)\cdot B_l^{(t)}(t)
\end{equation}
with $K_s = K_t = 12$ basis functions per dimension (\texttt{k = c(12,12)})
and \texttt{d = c(2,1)} grouping the two spatial coordinates into one
isotropic block.
The smoothness penalty acts as a Gaussian prior on the coefficients
$\boldsymbol{\beta}^{(k)}$, regularising the fit against overfitting
(\citealt{Nalborczyk2025}, Section~1.4).

\paragraph{MCMC settings.}
Two chains of 1500 iterations each (750 warmup) are run via
\texttt{cmdstanr} with four within-chain threads, yielding
$S = 1500$ posterior draws.
Because the posterior odds are bounded by $S$:
\begin{equation}
  \widehat{\mathrm{OR}}_\text{max} = S = 1500,
  \qquad
  \widehat{\mathrm{OR}}_\text{min} = 1/S \approx 6.7 \times 10^{-4}
\end{equation}
this is less precise than the 24\,000-draw setup of \citet{Nalborczyk2025}
but sufficient for pilot validation.

% ────────────────────────────────────────────────────────────────────────────
\subsection{Validation Targets}\label{sec:targets}

Three binary pass/fail criteria are evaluated after fitting.

\subsubsection*{Target 1 --- Convergence}

Let $N_\text{div}$ be the number of divergent HMC transitions,
$\hat{R}_\text{max}$ the largest Gelman--Rubin statistic across all
parameters \citep{Vehtari2021}, and $\text{ESS}_\text{min}$ the smallest
bulk effective sample size.
\begin{equation}\label{eq:t1}
  \text{T1} = (N_\text{div} = 0)
              \;\wedge\; (\hat{R}_\text{max} < 1.01)
              \;\wedge\; (\text{ESS}_\text{min} > 200)
\end{equation}
This mirrors the diagnostics reported in \citet{Nalborczyk2025}
(Section~2.1.2, p.\,9).

\subsubsection*{Target 2 --- Onset and offset recovery at CZ}

From $S^* = 1000$ posterior predictive draws at the CZ location, compute
the condition difference at each retained time point $t \in \mathcal{T}^*$:
\begin{equation}
  \Delta^{(s)}(t)
    = \hat{y}^{(s)}_{\text{CZ},t,k=2}
    - \hat{y}^{(s)}_{\text{CZ},t,k=1}
\end{equation}
Then form the posterior odds (cf.\ Equations~\ref{eq:prob}--\ref{eq:odds}):
\begin{equation}
  \widehat{\mathrm{OR}}(t)
    = \frac{\hat{p}(t)}{1 - \hat{p}(t)},
  \qquad
  \hat{p}(t) = \frac{1}{S^*}\sum_{s=1}^{S^*}\mathbf{1}[\Delta^{(s)}(t)>0]
\end{equation}
Onset and offset estimates follow \citet{Nalborczyk2025} (Figure~4):
\begin{equation}
  \hat{t}_\text{onset}  = \min\bigl\{t : \widehat{\mathrm{OR}}(t) > \theta\bigr\},
  \qquad
  \hat{t}_\text{offset} = \max\bigl\{t : \widehat{\mathrm{OR}}(t) > \theta\bigr\}
\end{equation}
with $\theta = 10$.
The pass criterion (20\,ms tolerance) is:
\begin{equation}\label{eq:t2}
  \text{T2} =
    \bigl|\hat{t}_\text{onset}  - 0.160\bigr| \le 0.020
  \;\wedge\;
  \bigl|\hat{t}_\text{offset} - 0.342\bigr| \le 0.020
  \quad (\text{seconds})
\end{equation}

\subsubsection*{Target 3 --- Spatial localisation}

At the midpoint time $t^* \approx (0.160 + 0.342)/2 = 0.251$\,s
(nearest grid point), compute the posterior mean condition difference
per electrode:
\begin{equation}
  \bar{\Delta}_c
    = \frac{1}{S^*}\sum_{s=1}^{S^*}
      \bigl(\hat{y}^{(s)}_{c,t^*,k=2} - \hat{y}^{(s)}_{c,t^*,k=1}\bigr)
\end{equation}
The pass criterion is that the top-ranked electrode lies in the immediate
neighbourhood of the true activation centre CZ:
\begin{equation}\label{eq:t3}
  \text{T3} = \Bigl(\arg\max_c\,\bar{\Delta}_c\Bigr)
              \in \{\text{CZ, FCZ, CPZ, C1, C2}\}
\end{equation}

\subsubsection*{Target 4 --- Peak amplitude recovery at CZ}

At the true peak time $t^* \approx 0.251$\,s, the posterior mean
condition difference at CZ should approximate the simulated peak
amplitude $A = 1.5\,\mu\mathrm{V}$. The pass criterion requires the
95\% credible interval to cover the truth:
\begin{equation}\label{eq:t4}
  \text{T4} = \bigl(\bar{\Delta}_{\text{CZ}, t^*}^{2.5\%}
                    \le A\bigr)
            \;\wedge\;
            \bigl(\bar{\Delta}_{\text{CZ}, t^*}^{97.5\%}
                    \ge A\bigr).
\end{equation}
Target~4 is sensitive to kernel-smoothing bias at the peak, where the
second derivative is maximally negative
(Section~\ref{sec:amp_theory}).

\subsubsection*{Target 4b --- Mean amplitude recovery over active window}

Recognising that ERP analyses commonly prefer mean amplitude over a
time window to peak amplitude \citep[Ch.~6]{Luck2014}, we complement
Target~4 with a window-averaged metric. Let $\mathcal{W} = \{t :
t_\text{onset} \le t \le t_\text{offset}\} = [0.160, 0.342]$\,s. For each
posterior draw $s$, define
\begin{equation}
  \bar{\mu}^{(s)}_{\text{CZ}, \mathcal{W}}
    = \frac{1}{|\mathcal{W}|}
      \sum_{t \in \mathcal{W}}
      \bigl(\hat{y}^{(s)}_{\text{CZ},t,k=2}
           - \hat{y}^{(s)}_{\text{CZ},t,k=1}\bigr),
\end{equation}
with sample-grid average. The truth is the corresponding window-mean
of the normalised Gaussian template times $A$:
\begin{equation}
  \mu^*_\mathcal{W}
    = A \cdot \frac{1}{|\mathcal{W}|}
      \int_\mathcal{W} f(t)\, dt
    \approx 0.731\,\mu\mathrm{V}.
\end{equation}
The pass criterion mirrors Target~4 (95\% CrI coverage):
\begin{equation}\label{eq:t4b}
  \text{T4b} = \bigl(\bar{\mu}^{2.5\%} \le \mu^*_\mathcal{W}\bigr)
            \;\wedge\;
            \bigl(\bar{\mu}^{97.5\%} \ge \mu^*_\mathcal{W}\bigr).
\end{equation}
Target~4b is theoretically less biased than Target~4: the integral
\(\int_\mathcal{W} f''(t)\,dt\) is smaller than $f''(t^*)$ because
$f''$ changes sign within the window (negative near peak, positive in
tails). Comparing Target~4 and Target~4b empirically quantifies the
peak-vs-mean bias gap predicted by classical kernel-smoothing theory.

Table~\ref{tab:t4_comparison} reports both metrics across our pilot
fits.

\begin{table}[ht]
\centering
\caption{Comparison of peak (Target~4) and window-mean (Target~4b)
amplitude recovery across the pilot BGAM fits. Truth peak
$A = 1.5\,\mu$V; truth mean over $\mathcal{W} = [160, 342]$\,ms is
$\approx 0.731\,\mu$V.}
\label{tab:t4_comparison}
\small
\begin{tabular}{lcccccccc}
\toprule
Fit & $k_{xy}$ & $k_t$ & full & $P$ &
  T4 recovery & T4b recovery & $\Delta$ (pp) \\
\midrule
\texttt{v4\_k12}        & 10 & 12 & --- & 1 & 37\%        & 55\%        & +18 \\
\texttt{v4\_k20}        & 10 & 20 & --- & 1 & 46\%        & 55\%        & +9  \\
\texttt{v5\_k20}        & 20 & 20 & --- & 5 & 46\%        & \textbf{76\%} & \textbf{+30} \\
\texttt{v5\_k3020\_full}& 30 & 20 & TRUE & 5 & 36\%        & 58\%        & +22 \\
\bottomrule
\end{tabular}
\end{table}

Across all four configurations, replacing peak with window-mean as
the amplitude target improves recovery by 9--30 percentage points,
qualitatively confirming the kernel-smoothing bias argument of
Section~\ref{sec:amp_theory}. The recommended configuration
\texttt{v5\_k20} achieves the largest improvement (+30\,pp) and the
best mean-amplitude recovery (76\%), reinforcing the choice
$(k_{xy}, k_t) = (20, 20)$.

\subsubsection*{Joint verdict}

\begin{equation}\label{eq:pilot_pass}
  \text{Pilot PASS} = \text{T1} \wedge \text{T2} \wedge \text{T3}
\end{equation}

Table~\ref{tab:targets} summarises the correspondence between these
targets and the evaluation framework of \citet{Nalborczyk2025}.

\begin{table}[ht]
\centering
\caption{Validation targets and their correspondence to
\citet{Nalborczyk2025}.}
\label{tab:targets}
\small
\begin{tabular}{p{1cm} p{4.5cm} p{7cm}}
\toprule
\textbf{Target} & \textbf{Criterion} & \textbf{Nalborczyk \& Bürkner (2025)} \\
\midrule
T1 & $\hat{R} < 1.01$, ESS $> 200$
   & Section~2.1.2: ``convergence assessed via trace plots
     and $\hat{R}$'' \\[4pt]
T2 & $|\hat{t} - t^*| \le 20$\,ms at CZ
   & Table~1: BGAM bias $\approx 0.1$\,ms onset,
     $2.4$\,ms offset; Figure~4: $\theta = 10$ \\[4pt]
T3 & Top electrode $\in$ CZ neighbourhood
   & Appendix~A, Figure~A2: spatial cluster
     centred on true source \\
\bottomrule
\end{tabular}
\end{table}

% ════════════════════════════════════════════════════════════════════════════
\section{Evaluation Pipeline (\texttt{03\_evaluate\_and\_plot.R})}\label{sec:eval}

This section describes how the fitted model is evaluated against ground truth,
following the framework of \citet{Nalborczyk2025} (Sections~2.1.3--2.1.4,
Table~1) extended to the spatial dimension.

% ────────────────────────────────────────────────────────────────────────────
\subsection{Spatial Weight Reconstruction}\label{sec:weight_recon}

Before evaluation, the true signal weight of each electrode is recomputed
from the data-generation model (Equation~\ref{eq:weights}):
\begin{equation}
  w_c = A\cdot\exp\!\left(-\frac{d_c^2}{2\sigma_s^2}\right),
  \qquad d_c = \|\hat{\mathbf{x}}_c - \hat{\mathbf{x}}_\text{CZ}\|_2
\end{equation}
with $A = 1.5$ and $\sigma_s = 1.2$ (standardised units).
These weights encode how much true signal each electrode carries and serve
as the basis for all evaluation metrics below.

% ────────────────────────────────────────────────────────────────────────────
\subsection{Binary Evaluation Metrics}\label{sec:binary}

The binary active set is defined via a distance cutoff equivalent to
$w_c > 0.1$:
\begin{equation}
  \mathcal{C}_\text{active} = \{c : d_c < 2.15\,\sigma_s\}
\end{equation}
For a detected electrode set $\hat{\mathcal{C}}$, let
$\mathrm{TP} = |\hat{\mathcal{C}} \cap \mathcal{C}_\text{active}|$,
$\mathrm{FP} = |\hat{\mathcal{C}} \setminus \mathcal{C}_\text{active}|$,
$\mathrm{FN} = |\mathcal{C}_\text{active} \setminus \hat{\mathcal{C}}|$.
The standard spatial metrics are:
\begin{equation}\label{eq:binary_metrics}
  \text{Precision} = \frac{\mathrm{TP}}{\mathrm{TP}+\mathrm{FP}}, \qquad
  \text{Recall}    = \frac{\mathrm{TP}}{\mathrm{TP}+\mathrm{FN}}, \qquad
  F_1 = \frac{2\cdot\text{Precision}\cdot\text{Recall}}
             {\text{Precision}+\text{Recall}}
\end{equation}

% ────────────────────────────────────────────────────────────────────────────
\subsection{Gaussian-Weighted Evaluation}\label{sec:weighted}

A hard threshold treats $w_c = 0.09$ identically to $w_c \approx 0$,
discarding meaningful spatial gradient information.
The weighted metrics instead use $w_c$ directly.

Define the normalised weight distribution
$\tilde{w}_c = w_c / \sum_{c'} w_{c'}$.
The \textbf{weighted recall} is the fraction of total signal mass detected:
\begin{equation}\label{eq:w_recall}
  \mathrm{Recall}_w = \sum_{c\,\in\,\hat{\mathcal{C}}} \tilde{w}_c
\end{equation}
The \textbf{weighted precision} is the average weight of detected electrodes
relative to the overall mean:
\begin{equation}\label{eq:w_prec}
  \mathrm{Precision}_w
  = \frac{\bar{w}_{\hat{\mathcal{C}}}}{\bar{w}_\text{all}}
  = \frac{\frac{1}{|\hat{\mathcal{C}}|}\sum_{c\,\in\,\hat{\mathcal{C}}} w_c}
         {\frac{1}{C}\sum_{c=1}^{C} w_c}
\end{equation}
$\mathrm{Precision}_w > 1$ means detected electrodes are, on average,
closer to the activation centre than chance---a desirable property.
The weighted $F_1$ is:
\begin{equation}\label{eq:w_f1}
  F_{1,w} =
  \frac{2\cdot\mathrm{Precision}_w\cdot\mathrm{Recall}_w}
       {\mathrm{Precision}_w + \mathrm{Recall}_w}
\end{equation}
No threshold parameter is required; the evaluation depends only on the
Gaussian decay structure of the true signal.

% ────────────────────────────────────────────────────────────────────────────
\subsection{\boldmath$k\sigma$ Robustness Check}\label{sec:ksigma}

Because the binary cutoff $k_0 = 2.15$ is a modelling choice, we vary $k$
and recompute all binary metrics with the detection result $\hat{\mathcal{C}}$
held fixed:
\begin{equation}\label{eq:ksigma}
  \mathcal{C}_\text{active}(k) = \{c : d_c \leq k\,\sigma_s\},
  \qquad k \in \{1.0,\;1.5,\;2.0,\;2.15,\;2.5,\;3.0\}
\end{equation}
A \textbf{flat curve} across $k$ indicates robust conclusions;
a \textbf{steep curve} signals sensitivity to the ground-truth definition.
This is the spatial analogue of the threshold sensitivity in
\citet{Nalborczyk2025} (Section~2.1.3): in both cases the expensive
model fit is reused and only the evaluation criterion changes.

% ────────────────────────────────────────────────────────────────────────────
\subsection{Posterior Odds Threshold Sensitivity}\label{sec:theta_sens}

Five thresholds are evaluated without refitting the model:
\begin{equation}
  \theta \in \{3,\;5,\;10,\;19,\;99\}
\end{equation}
For each $\theta$, the active spatio-temporal node set
$\mathcal{A}(\theta) = \{(c,t): \widehat{\mathrm{OR}}_{c,t} > \theta\}$
is computed, spatial clusters are extracted via connected components of the
electrode adjacency graph, and onset bias, offset bias, Precision, Recall,
and cluster count are recorded for the largest cluster.

% ────────────────────────────────────────────────────────────────────────────
\subsection{Topographic Visualisation}\label{sec:topo}

At $N = 8$ snapshot times $t_1^*,\ldots,t_8^*$ evenly spaced across the
trial, a thin-plate spline GAM ($k = 40$ basis functions) is fitted to
the $C$ electrode values to produce a continuous scalp surface
$\hat{z}(x,y)$ on an $80\times 80$ grid.
Four rows of panels are plotted:

\begin{equation}\label{eq:topo_rows}
\begin{aligned}
  \text{Row 1 (GT signal):}   &\quad
    s_{c,t^*} = A\cdot w_c \cdot f(t^*;\,\overline{\mathrm{onset}})\\
  \text{Row 2 (observed):}    &\quad
    \bar{y}_{c,t^*,2} - \bar{y}_{c,t^*,1}\\
  \text{Row 3 (BGAM predicted):} &\quad
    \hat{f}_\Delta(\hat{x}_c,\hat{y}_c,t^*)\\
  \text{Row 4 (residual):}    &\quad
    (\bar{y}_{c,t^*,2} - \bar{y}_{c,t^*,1})
    - \hat{f}_\Delta(\hat{x}_c,\hat{y}_c,t^*)
\end{aligned}
\end{equation}

Row~1 uses the noiseless ground-truth signal evaluated at the mean onset
$\overline{\mathrm{onset}} = 160$\,ms, providing a reference against which
Rows~2--4 can be directly compared.

\paragraph{Note on current code (missing Row 1).}
The current \texttt{03\_evaluate\_and\_plot.R} produces only Rows~2--4.
Row~1 is absent because the ground truth signal surface has not been
computed from the stored \texttt{ground\_truth\$sigma\_space},
\texttt{ground\_truth\$bump\_amplitude}, and the electrode coordinates.
This should be added as:
\begin{verbatim}
gt_signal <- sensors_tbl %>%
    mutate(signal = ground_truth$bump_amplitude * weight *
                    erp_at(t_star, onset = 0.160))
\end{verbatim}
Without Row~1, it is impossible to visually verify whether the Gaussian
spatial bump is centred correctly and has the expected width---a key
sanity check before interpreting the BGAM results.

% ════════════════════════════════════════════════════════════════════════════
\section{Theoretical limits on amplitude recovery}\label{sec:amp_theory}

Across our diagnostic suite (Section~\ref{sec:targets}), peak amplitude at the
seed electrode (Target~4) is systematically under-recovered, even when
convergence (Target~1), spatial localisation (Target~3), and temporal-window
detection (Target~2) succeed. This section explains why this dissociation is
expected from classical nonparametric-regression theory and is, in fact, the
canonical fingerprint of smoothing-induced bias.

\subsection{Kernel-smoothing bias at peaks}\label{sec:bias_formula}

For any smoothing-spline or kernel estimator with effective bandwidth $h$,
the pointwise expected bias has the asymptotic expansion
\begin{equation}\label{eq:wand_jones}
\mathbb{E}[\hat{f}(x_0)] - f^{*}(x_0)
   \;=\; \tfrac{1}{2}\, h^{2}\, f^{*\prime\prime}(x_0)\, \mu_{2}(K)
        + O(h^{4}),
\end{equation}
where $f^{*}$ is the true regression function, $\mu_{2}(K)=\int u^{2}K(u)\,du$
is the second moment of the equivalent kernel, and $h$ depends monotonically
on the smoothing parameter $\lambda$ \citep[\S 4.3]{WandJones1995}; see also
\citet[\S 3.1]{FanGijbels1996}. Because smoothing splines are mathematically
equivalent to kernel smoothers with a particular kernel \citep{Eubank1988,
GreenSilverman1994}, Equation~\eqref{eq:wand_jones} applies directly to
GAM/BGAM estimates.

At a local maximum, $f^{*\prime\prime}(x_0) < 0$, so the bias is negative:
the estimator systematically \emph{under-estimates} peak amplitude. The
magnitude scales with the curvature of the true function and with $h^{2}$.

\paragraph{Specialisation to a Gaussian peak.}
For a truncated-Gaussian ERP with amplitude $A$ and temporal scale
$\sigma_{\mathrm{ERP}}$, the second derivative at the peak is
$f^{*\prime\prime}(t_0) = -A/\sigma_{\mathrm{ERP}}^{2}$, so
\begin{equation}\label{eq:bias_gauss}
\mathbb{E}[\hat{f}(t_0)] - A
   \;\approx\; -\, \frac{A\, h^{2}\, \mu_{2}(K)}{2\,\sigma_{\mathrm{ERP}}^{2}}.
\end{equation}
With our pilot parameters ($A = 1.5\,\mu\mathrm{V}$,
$\sigma_{\mathrm{ERP}} = 50\,\mathrm{ms}$) and an effective bandwidth on
the order of the data-adaptive smoothing scale
($h \sim 50\,\mathrm{ms}$, $\mu_{2}(K)\approx 1$),
Equation~\eqref{eq:bias_gauss} predicts a bias of roughly $-0.75\,\mu\mathrm{V}$,
corresponding to a point-estimate recovery near 50\%. This matches the
36--46\% recovery observed across our 3D HMC fits
(\texttt{v5\_k20}, \texttt{v5\_k3020\_full}), well within the precision of the
$O(h^{4})$ asymptotic approximation. A linear local-polynomial smoother (which
includes our $m=2$ thin plate spline) cannot remove this second-order curvature
bias \citep[\S 3.1]{FanGijbels1996}; quadratic local polynomials are the
lowest-order option able to correct it.

\subsection{Multivariate extension and the curse of dimensionality}\label{sec:bias_mvt}

For a $d$-dimensional smoother with marginal bandwidths $h_{1},\dots,h_{d}$,
Equation~\eqref{eq:wand_jones} generalises to
\begin{equation}\label{eq:bias_mvt}
\mathbb{E}[\hat{f}(\mathbf{x}_0)] - f^{*}(\mathbf{x}_0)
   \;\approx\; \tfrac{1}{2}\,\sum_{j=1}^{d}\, h_{j}^{2}\,
   \frac{\partial^{2} f^{*}}{\partial x_{j}^{2}}(\mathbf{x}_0)\, \mu_{2}(K_{j}).
\end{equation}
Our spatio-temporal smooth has $d = 3$ (two spatial dimensions plus time).
At the seed electrode and peak time, all three second derivatives are
negative: the true signal is a local maximum in space (Gaussian spatial decay
centred on CZ, Equation~\eqref{eq:spatial_weight}) and in time (Gaussian ERP,
Section~\ref{sec:erp}). Consequently the bias terms accumulate rather than
cancel.

Two factors amplify the bias further in our 3D setting relative to a 1D
fit:
\begin{enumerate}
  \item \textbf{Information density.} With $P=5$ participants and
        $R=15$ trials, the 1D fit has roughly 510 cells for $\sim$40 smooth
        coefficients (13 cells per parameter), while the 3D fit
        ($\verb|t2(xproj, yproj, time)|$, $k=c(30,20)$, \texttt{full=TRUE},
        \verb|by=condition|) has $\sim$5100 cells for $\sim$1800 coefficients
        (2.8 cells per parameter). The likelihood is therefore far less
        informative in 3D, and the posterior is closer to the prior. Classical
        Bayesian shrinkage theory \citep{LindleySmith1972,Efron1971} predicts
        a larger posterior-mean attenuation toward zero.
  \item \textbf{Inflated measurement-error term.}
        \texttt{aggregate\_trials()} reports \texttt{eeg\_sd}, the
        \emph{trial-level} standard deviation. Passing this to
        \verb|se(eeg_sd)| overstates the standard error of \texttt{eeg\_mean}
        by a factor of $\sqrt{R}\approx 3.87$ (the correct quantity is
        $\texttt{eeg\_sd}/\sqrt{R} = \texttt{eeg\_sem}$). The effective
        bandwidth chosen by the model is consequently inflated, propagating
        through Equation~\eqref{eq:bias_mvt} as a multiplicative penalty on
        peak recovery.
\end{enumerate}

These factors are individually mild and almost invisible in the original 1D
fit of \citet{Nalborczyk2025} on dense data (97.5\% recovery for
$P=20, R=50$ even with \verb|se(eeg_sd)|), but combine destructively in the
sparse 3D setting of our pilot.

\subsection{Implications and mitigations}\label{sec:amp_mitigation}

The bias is a property of the estimator under the chosen specification, not
a numerical bug. Strategies that may reduce it---each with trade-offs---are:
\begin{itemize}
  \item Use the standard-error of the mean, $\verb|se(eeg_sd/sqrt(R))|$ or
        the precomputed \texttt{eeg\_sem}, when fitting on aggregated cells
        (Section~\ref{sec:reduction}).
  \item Loosen the prior on the smoothing-variance components
        (\texttt{class = "sds"}), so that the likelihood---rather than the
        prior---selects the effective bandwidth.
  \item Use higher-order local-polynomial smoothers (e.g.\
        $m = 3$ thin plate splines) which can remove the leading
        second-derivative bias \citep[\S 3.1]{FanGijbels1996}.
  \item Increase data density by using the full $61$-electrode layout and
        $R = 50$ trials per condition (Table~\ref{tbl:fixed}); this directly
        reduces the curse-of-dimensionality contribution.
  \item Report \emph{both} amplitude and timing diagnostics; do not rely on
        BGAM amplitudes for effect-size inference in the sparse-data regime
        without an explicit calibration analysis (analogous to the general
        recommendation of \citet[Ch.~6]{Luck2014} to prefer mean over peak
        amplitude in noisy ERPs).
\end{itemize}

\paragraph{Relation to existing literature.}
The classical bias formulae above are textbook material
\citep{WandJones1995,FanGijbels1996}, and the over-shrinkage of interaction
components in tensor-product ANOVA decompositions is documented in
\citet[\S 3]{Gu2002} and inherited by the \texttt{t2()} construction
\citep{WoodScheiplFaraway2013}. Related phenomena in the EEG literature
include latency-jitter-induced amplitude attenuation in trial-averaged
ERPs \citep[Ch.~6]{Luck2014} and filter-induced amplitude bias
\citep{Acunzo2012}. However, to our knowledge, no prior published work has
systematically benchmarked peak-amplitude recovery in 3D Bayesian GAMs
applied to simulated ERPs. The present results therefore fill a documented
gap between general nonparametric theory and the specific BGAM-EEG
application proposed by \citet{Nalborczyk2025}.

% ════════════════════════════════════════════════════════════════════════════
\bibliography{refs}

\end{document}
